\section{Revisión a nivel de sistema: qué combina realmente GLOM}

Llegados a este punto, GLOM puede condensarse en un recurrent structured inference loop: la evidence de bajo nivel propone hacia arriba una parent hypothesis, el context de alto nivel predice hacia abajo la part, la same-level attention hace que las hypothesis similares formen islands, el previous state aporta inertia, la masked reconstruction y la consensus distillation entrenan estas funciones, y el location-conditioned decoder se encarga de generar, a partir de un shared whole, las distintas parts.

\lecturefigure{slide-32-summary.jpg}{Síntesis de la clase: combinar Transformer, contrastive learning y neural field para construir GLOM.}{Video oficial de Stanford Online 00:50:08--00:52:20.}

\subsection{Lista de mecanismos y modos de falla}

\begin{table}[H]
\centering
\small
\caption{Mecanismos de GLOM, problemas que resuelven y evidencia que todavía falta.}
\begin{tabularx}{\textwidth}{L{0.20\textwidth}L{0.30\textwidth}>{\raggedright\arraybackslash}X}
\toprule
Mecanismo & Problema que intenta resolver & Validación pendiente \\
\midrule
Dynamic embeddings & Representar con hardware fijo el parse específico de cada entrada & Alignment con la ground-truth hierarchy \\
Islands of agreement & Representar un mismo node con la misma activity & Evitar el collapse, separación entre instancias, estabilidad de las fronteras \\
Bottom-up/top-down maps & Consistencia de identity/pose entre part y whole & Coordinate equivariance y occlusion robustness \\
Local similarity attention & Descubrir regiones de spatial coherence & Attractor, error propagation y convergence \\
Hough-style parent voting & Evitar la explosión del pairwise relation routing & Desempeño en scenes complejas, con múltiples instancias y con deformable objects \\
Consensus distillation & Hacer que los local predictors compartan información estructural & Teacher drift, consensos erróneos y optimization stability \\
Sparse upper levels & Ampliar el range sin disparar el cost & FLOPs, memory, latency y communication reales \\
Neural field decoder & Que un mismo whole prediga parts distintas en posiciones distintas & Calidad de generación tras el entrenamiento conjunto con las islands \\
\bottomrule
\end{tabularx}
\end{table}

\subsection{Experimentos que conviene exigir al leer el artículo}

Si hoy se quisiera llevar GLOM de un design document a un artículo de sistemas, haría falta, como mínimo, cerrar el ciclo con los siguientes experimentos:

\begin{enumerate}
  \item \textbf{Representation readout}: definir cómo se extraen nodes, edges, identity y pose a partir de los embeddings continuos;
  \item \textbf{Causal use}: tras intervenir una island, verificar si la prediction de la part/whole correspondiente cambia conforme a la estructura;
  \item \textbf{Convergence}: medir settling steps, fixed-point stability, oscillation y early stopping;
  \item \textbf{Anti-collapse}: comparar local attention, negative pairs, variance regularization y consensus weighting;
  \item \textbf{Generalization}: evaluar con viewpoints nuevos, part compositions nuevas, occlusion, multiple instances y video tracking;
  \item \textbf{Systems accounting}: reportar state memory, attention edges, BPTT activation, wall-clock latency y energy;
  \item \textbf{Ablation}: retirar por separado top-down, bottom-up, persistence, lateral attention, location condition y distillation.
\end{enumerate}

\begin{importantbox}{Una afirmación falsable sobre GLOM}
Si las islands realmente desempeñan el papel de parse-tree node, retirar o perturbar una island del parent-level debería alterar de forma sistemática la top-down prediction de sus constituent parts, y no solo reducir la puntuación de clasificación global. Esta intervention se acerca más a una evidencia del mecanismo que el argumento de que «la visualización parece un object».
\end{importantbox}

\subsection{Resumen de la sección}

La fortaleza de GLOM es reunir representation, inference, learning y coordinate generation en un mismo lenguaje de diseño; su debilidad es que casi todos sus eslabones clave carecen de experimentos de extremo a extremo. Precisamente porque la evidencia está incompleta, resulta especialmente adecuado como architecture thinking exercise: obliga al lector a reescribir la «estructura de los objetos», que era un término intuitivo, en forma de state, message, loss y verification contract.
